\scanpage{23}
\mspage{21}
$\operatorname{DR}^1(A/S)$ est un foncteur \deleted{\unsure{additi|additif}} multiplicatif en $A$, grâce à Kunneth postulé dans 2.2., et que l'on a $FV=VF=p.\operatorname{id}$ sur les schémas abéliens en car.$p$, on en conclut les mêmes relations dans $\operatorname{DR}^1$. Cela montre donc que $\operatorname{DR}^1(A)$ est un bic\deleted{\unsure{c}}ristal de poids 1, i.e. un \uline{cristal de Dieudonné}. Lorsque $S$ est le spectre d'un corps \inserted{parfait} de caractéristique $p>0$, un tel cristal s'identifie simplement, d'après 1.12, à un module de Dieudonné sur $W=W(k)$. Bien sûr, \uline{on doit trouver que ce module n'est autre que le module de Dieudonné du groupe formel} \deleted{défini} (ou plutôt, du groupe p-divisi\unsure{ble|xble}) \uline{défini par la variété abélienne} $A$. Débarassé de l'hyperstructure axiomatico-conjectural de 2.2., \deleted{en fait} on peut exprimer ainsi la construction \deleted{ainsi} obtenue du module de Dieudonné:

1º) Soit $A$ un schéma abélien sur un schéma affine $S$. On sait que pour tout morphisme $S\longrightarrow T$ d'immersion fermée surjective, défini par Idéal nilpotent sur $T$, $A$ se prolonge en un schéma abélien $B$ sur $T$. On peut regarder la cohomologie de De-Rham ordinaire $\underline{R}^1g_*(\underline{\Omega}^{\bullet}_{B/T})$, où $g:B\longrightarrow T$ est le morphisme structural, et on sait que c'est un Module localement libre de rang $2g$, où $g$ est la dimension relative de $A$ sur $S$ donc de $B$ sur $T$. \uline{Première affirmation}: ce Module, à isomorphisme canonique près, ne dépend que \inserted{de $A$ sur $S$ et} de $S\longrightarrow T$, et pas du choix du prolongement $B$ de $A$. En d'autres termes, la donnée de $A$ sur $S$ définit un \uline{cristal} localement libre de rang $2g$ sur $S$, soit $\operatorname{DR}^1(A/S)$. C'est évidemment un foncteur additif en $A$, compatible avec le changement de base, \inserted{et définissable sans condition \unsure{affine} sur $S$ (en recollant).}

2º) Supposons que $S$ soit le spectre d'un corps parfait $k$ de car $p>0$, alors le cristal précédent s'identifie à un module libre de rang $2g$ sur $W=W(k)$. On y introduit les structures $F$ et $V$, en utilisant comme ci-dessus les homomorphismes $F:A\longrightarrow A^{(p)}$ et $V:A^{(p)}\longrightarrow A$. On trouve ainsi un module de Dieudonné \inserted{$M$}, et: \uline{Deuxième affirmation}: C'est bien celui défini par Dieudonné. En d'autres termes: si $\Lambda$ est n'importe quel anneau local artinien de corps résiduel $k$, et $B$ un prolongement de $A$ en un schéma abélien sur $\Lambda$, alors \deleted{\unsure{$H$|la}} la cohomologie de De Rham $H^1_{\operatorname{DR}}(B)=\underline{H}^1(B,\underline{\Omega}^{\bullet}_{B/\Lambda})$ est canoniquement \inserted{et fonctoriellement} isomorphe à $M\otimes_W\Lambda$, où $M$ est le module de Dieudonné classique de $A$, et où on tient compte du th. de Cohen qui munit $\Lambda$ d'une structure canonique de $W$-algèbre. (NB la \deleted{canonicité} fonctorialité de
\margin{22}
